% CLOSE -- os/samir/samir_main.c repl_do_close (Using III+ U5-53/PDF 165:
% CLOSE <file type> / ALL.  ALL and DATABASES close every database with its
% indexes; INDEX closes the index files of every work area and leaves the
% databases open in natural order; ALTERNATE, FORMAT and PROCEDURE are
% accepted and close nothing -- SET ALTERNATE TO closes the alternate file,
% observed in the replay; any other word, or none, is 10).  Lane samirdml 14a8fc75.
% Gates: test-samir-repl, test-samir-dml, test-samir-dml-emu.
\begin{XbaseEntry}{CLOSE}{command}
Closes the open files of one type, or all of them.
\Syntax
\SyntaxLine{\kw{CLOSE} \xp{file type}}
\par or\par
\SyntaxLine{\kw{CLOSE ALL}}
\Usage
\cmd{CLOSE} is followed by one of the words \kw{ALL}, \kw{DATABASES},
\kw{INDEX}, \kw{ALTERNATE}, \kw{FORMAT} or \kw{PROCEDURE}.  Memory variables
are not affected by any of them.

\cmd{CLOSE ALL} and \cmd{CLOSE DATABASES} close the database open in every
work area, with its index files.  Afterwards no database is in USE in any
work area.

\cmd{CLOSE INDEX} closes the index files of every work area and leaves the
databases open.  Each database is then read in the order of its records
until an index file is opened again with \cmd{SET INDEX} or \cmd{USE}.

\cmd{CLOSE ALTERNATE}, \cmd{CLOSE FORMAT} and \cmd{CLOSE PROCEDURE} are
accepted and close nothing in this version.  Close the alternate file with
\cmd{SET ALTERNATE TO}, without a file name.  \cmd{CLOSE} without a word, or with any other word,
is reported as \xmsg{10}{Syntax error.}
\Tips
To close only the database of the selected work area, type \cmd{USE} with no
file name.
\Examples
\DotExample{examples/c-close}
With its index file closed, the database is read in the order of its
records again:
\DotExample{examples/c-close-2}
\SeeAlso{QUIT, SELECT, USE, SET INDEX}
\end{XbaseEntry}

